% ---------------------------------------------------------------------------
% RoboForce - Robotics Electrical Engineering Intern
% Tailored from bases/embedded_controls.tex on 2026-08-25
% Worksheet: notes.md in this folder. Posting: job_description.txt
%
% Do NOT put the company name anywhere in the rendered document. It belongs in
% the filename only.
% ---------------------------------------------------------------------------
\documentclass[10pt,letterpaper]{article}
\usepackage{resume}

% Keep adjacent roles compact while preserving clear visual separation.
\setlist[itemize]{topsep=1pt,itemsep=0.5pt,after=\vspace{1pt}}
\renewcommand{\entryspace}{\vspace{0pt}}

\begin{document}
\bodyfont

\resumeheader
\educationsection

\sectionheader{TECHNICAL SKILLS}

\techline{\textbf{Languages}: C/C++, Python, Assembly, Java}
\techline{\textbf{Embedded Firmware}: STM32H7, TTC 60 VCU, PX4, real-time control, watchdogs, microcontroller I/O}
\techline{\textbf{Protocols \& Linux}: CAN, SPI, Linux SocketCAN, python-can, WebSockets}
\techline{\textbf{Test \& Debug}: Board bring-up, software-in-the-loop, scripted replay, GDB, BusMaster, TTC Downloader, telemetry traces}

\sectionheader{EXPERIENCE}

\roleline{Software Lead}{Jun 2026 - Present}
\orgline{\spartanracing}
\begin{itemize}
  \item Leading an in-house VCU program replacing a vendor controller with custom STM32H755 hardware, running a \textbf{10 ms} deterministic control task under a watchdog that reboots the controller if the task hangs
  \item Built its dual-bus CAN driver and identified non-\textbf{5 V}-tolerant analog inputs during bring-up, specifying voltage dividers before vehicle integration
  \item Developing the custom-inverter path for a \textbf{four-motor} EV, mapping calibrated pedal input to bounded CAN current commands with fail-safe precharge and ready-to-drive gating
  \item Reworked custom-BMS integration, decoding \textbf{46 CAN message IDs} across \textbf{96 cells}, \textbf{96 thermistors}, and \textbf{8 modules}, with a \textbf{1 s} communication timeout
\end{itemize}

\entryspace

\roleline{Software Engineer Intern}{Feb 2026 - Present}
\orgline{Argus Defense \textbar{} San Jose, CA}
\begin{itemize}
  \item Developed and validated embedded flight-control firmware in C on PX4, using structured telemetry logs and repeatable flight tests to reproduce issues and evaluate revisions
  \item Integrated and tuned visual-inertial odometry for onboard state estimation, comparing configuration changes against logged position-hold behavior
\end{itemize}

\entryspace

\roleline{Software Controls Engineer}{Jul 2025 - Jun 2026}
\orgline{\spartanracing}
\begin{itemize}
  \item Developed fixed-point PI firmware with saturation, anti-windup, and slew limiting that held measured power within \textbf{0.01\%} of its commanded setpoint while enforcing an \textbf{80 kW} limit
  \item Built CAN diagnostics firmware for live power and torque frames, status bits, and error counters for real-time BusMaster monitoring and regression review
\end{itemize}

\sectionheader{PROJECT EXPERIENCE}

\roleline{STM32 Driver Display, Embedded Instrumentation}{Feb 2026 - Present}
\orgline{\spartanracing}
\techline{C, STM32H7, CAN, SPI}
\begin{itemize}
  \item Developed STM32H7 firmware that receives real-time CAN telemetry and drives an SPI LCD plus LED alerts for high-voltage, throttle, temperature, and power-budget status
\end{itemize}

\entryspace

\roleline{SR-Wireless-CAN, Linux Firmware Deployment \& Test Tool}{2026}
\orgline{\spartanracing \textbar{} Backend lead}
\techline{Python, Linux SocketCAN, python-can, FastAPI, WebSockets, SQLite, Docker}
\begin{itemize}
  \item Reverse-engineered and implemented an \textbf{8-stage} CAN firmware-deployment workflow from vendor traces, validating acknowledgements and CRC results over Linux SocketCAN
  \item Added simulated, recorded-trace, and live-CAN modes and documented frame encoding, timeouts, configuration, and recovery steps for repeatable debugging
\end{itemize}

\end{document}

